\section{Introduction}
\label{sec:introduction}

% TODO: write the introduction (≈1 to 1.5 pages).
% Suggested structure:
%   - Motivation: NER pipelines in production face an open-world setting - new
%     entity types appear, and re-annotating + retraining is expensive.
%   - Limitations of existing approaches:
%       * Closed-set NER (BIO tagging) cannot generalize beyond training labels.
%       * Triplet-loss + K-means (e.g. OWNER) requires training a dedicated encoder.
%       * Prompting LLMs is costly and brittle at scale.
%   - Our contribution: a lightweight open-world NER pipeline that reuses off-the-shelf
%     components (a frozen detector and a Matryoshka embedder) with only a light
%     contrastive fine-tuning plus hard-negative mining and a balanced linear typing
%     head; no encoder is trained from scratch. The pipeline is modular: each component
%     (detector, embedder, typing head) is independently swappable.
%   - Contributions list:
%       1. OPENER, an open-world NER pipeline assembled from off-the-shelf parts, with
%          no encoder trained from scratch.
%       2. Error-driven hard-negative mining that sharpens the embedding space for typing.
%       3. A zero-shot typing variant (OPENER-ZS) based on label-name prototypes.
%       4. An energy-aware benchmark of open-world NER on 13 datasets (AMI + latency +
%          energy), separating typing-on-gold from end-to-end.
%   - Paper roadmap (\S\ref{sec:related} discusses related work, \S\ref{sec:method}
%     details the OPENER pipeline, \S\ref{sec:experiments} evaluates on
%     [datasets], and \S\ref{sec:conclusion} concludes).

\emph{[Placeholder - introduction to be written.]}
